\documentclass[10pt,a4paper]{amsart}
\usepackage{amsmath,amssymb,amsthm,graphicx}
\usepackage{verbatim}
\newenvironment{reference}{\comment}{\endcomment}
%\newenvironment{reference}{}{}
\newenvironment{slogan}{\comment}{\endcomment}
\newenvironment{history}{\comment}{\endcomment}

% For commutative diagrams we use Xy-pic
\usepackage[all]{xy}

% We use 2cell for 2-commutative diagrams.
\xyoption{2cell}
% Theorem environments.
%
\theoremstyle{plain}
\newtheorem{theorem}[subsection]{Theorem}
\newtheorem{proposition}[subsection]{Proposition}
\newtheorem{lemma}[subsection]{Lemma}

\theoremstyle{definition}
\newtheorem{definition}[subsection]{Definition}
\newtheorem{example}[subsection]{Example}
\newtheorem{exercise}[subsection]{Exercise}
\newtheorem{situation}[subsection]{Situation}

\theoremstyle{remark}
\newtheorem{remark}[subsection]{Remark}
\newtheorem{remarks}[subsection]{Remarks}

\numberwithin{equation}{subsection}

% Macros
%
\def\lim{\mathop{\mathrm{lim}}\nolimits}
\def\colim{\mathop{\mathrm{colim}}\nolimits}
\def\Spec{\mathop{\mathrm{Spec}}}
\def\Hom{\mathop{\mathrm{Hom}}\nolimits}
\def\Ext{\mathop{\mathrm{Ext}}\nolimits}
\def\SheafHom{\mathop{\mathcal{H}\!\mathit{om}}\nolimits}
\def\SheafExt{\mathop{\mathcal{E}\!\mathit{xt}}\nolimits}
\def\Sch{\mathit{Sch}}
\def\Mor{\mathop{\mathrm{Mor}}\nolimits}
\def\Ob{\mathop{\mathrm{Ob}}\nolimits}
\def\Sh{\mathop{\mathit{Sh}}\nolimits}
\def\NL{\mathop{N\!L}\nolimits}
\def\CH{\mathop{\mathrm{CH}}\nolimits}
\def\proetale{{pro\text{-}\acute{e}tale}}
\def\etale{{\acute{e}tale}}
\def\QCoh{\mathit{QCoh}}
\def\Ker{\mathop{\mathrm{Ker}}}
\def\Im{\mathop{\mathrm{Im}}}
\def\Coker{\mathop{\mathrm{Coker}}}
\def\Coim{\mathop{\mathrm{Coim}}}

% Boxtimes
%
\DeclareMathSymbol{\boxtimes}{\mathbin}{AMSa}{"02}

%
% Macros for moduli stacks/spaces
%
\def\QCohstack{\mathcal{QC}\!\mathit{oh}}
\def\Cohstack{\mathcal{C}\!\mathit{oh}}
\def\Spacesstack{\mathcal{S}\!\mathit{paces}}
\def\Quotfunctor{\mathrm{Quot}}
\def\Hilbfunctor{\mathrm{Hilb}}
\def\Curvesstack{\mathcal{C}\!\mathit{urves}}
\def\Polarizedstack{\mathcal{P}\!\mathit{olarized}}
\def\Complexesstack{\mathcal{C}\!\mathit{omplexes}}
% \Pic is the operator that assigns to X its picard group, usage \Pic(X)
% \Picardstack_{X/B} denotes the Picard stack of X over B
% \Picardfunctor_{X/B} denotes the Picard functor of X over B
\def\Pic{\mathop{\mathrm{Pic}}\nolimits}
\def\Picardstack{\mathcal{P}\!\mathit{ic}}
\def\Picardfunctor{\mathrm{Pic}}
\def\Deformationcategory{\mathcal{D}\!\mathit{ef}}
\usepackage{cleveref}
\crefname{originaldefinition}{original definition}{original definitions}
\Crefname{originaldefinition}{Original definition}{Original definitions}
\makeatletter
\newlabel{definition-limit-spectral-sequence}{{011O}{source}{}{sourcepointer-1}{}}
\newlabel{definition-limit-spectral-sequence@cref}{{[originaldefinition][0][]011O}{[1][0][]source}{}{}{}}
\makeatother
\crefname{originallemma}{original lemma}{original lemmas}
\Crefname{originallemma}{Original lemma}{Original lemmas}
\makeatletter
\newlabel{lemma-spectral-sequence-filtered-differential}{{012C}{source}{}{sourcepointer-2}{}}
\newlabel{lemma-spectral-sequence-filtered-differential@cref}{{[originallemma][0][]012C}{[1][0][]source}{}{}{}}
\makeatother
\crefname{originaldefinition}{original definition}{original definitions}
\Crefname{originaldefinition}{Original definition}{Original definitions}
\makeatletter
\newlabel{definition-filtration-cohomology-filtered-differential}{{012E}{source}{}{sourcepointer-3}{}}
\newlabel{definition-filtration-cohomology-filtered-differential@cref}{{[originaldefinition][0][]012E}{[1][0][]source}{}{}{}}
\makeatother
\title{The limiting page of the spectral sequence and the associated graded of the cohomology}
\author{The Stacks Project Authors}
\date{stacks-project:homology.tex:local-snapshot-20261012; GNU FDL 1.2 or later}
\begin{document}
\maketitle
\noindent\textit{Editorial note.} All delivered mathematical text is the author's own, copied from the pinned Stacks Project chapter file homology.tex; no statement or proof step was rewritten, reordered or removed.\par\medskip
\noindent\textit{Editorial note.} The proof is the author's complete argument for both assertions; the parenthetical 'follows from an elementary argument about unions and intersections of graded subobjects' names the kind of argument used for a routine step whose inputs are displayed in the span, and it is preserved verbatim. No step of the proof is passed to the reader and the closing sentence is a conclusion.\par\medskip
\noindent\textit{Editorial note.} Every labelled reference occurring inside the counted spans is carried in assembly.external\_references as a non-load-bearing navigation pointer (same-chapter ones at their definition site, cross-chapter ones at the invocation inside the counted span), so no load-bearing mathematics is replaced by a stub. No bibliography citation occurs inside the counted spans.\par\medskip
\noindent\textit{Editorial note.} The same-name check for this card scanned every staged Stacks chapter of this repository -- all sources/stacks-*/pinned-original-source.tex files, the card's own chapter excluded -- and the labels found are recorded in dedup.same\_name\_cross\_chapter on the card.\par\medskip
\noindent\textit{Editorial note.} The preamble is the chapter's own native preamble excerpt already used for this source: the comment-type reference, slogan and history environments (native preamble lines 8-13), the Xy-pic setup (native preamble lines 14-18) and the native theorem declarations with the mathematical macros (native preamble lines 154-217, the actual last line of the file).\par\medskip
\noindent\textit{Allowed standard prerequisites (editorial).}\begin{itemize}
\item The limit of a spectral sequence. Same chapter, Definition definition-limit-spectral-sequence. Supplies the objects Z\_inf, B\_inf and the limiting object E\_inf = Z\_inf/B\_inf whose existence in each degree is the hypothesis of the statement.
\item The spectral sequence of a filtered differential object. Same chapter, Lemma lemma-spectral-sequence-filtered-differential. Supplies the construction of the spectral sequence together with the formulae for the boundaries and cocycles, which are what identify the degree p parts of Z\_inf and B\_inf.
\item The filtration on the cohomology of a filtered differential object. Same chapter, Definition definition-filtration-cohomology-filtered-differential. Supplies the displayed description of gr\textasciicircum{}p H(K) that the subquotient computed in the proof is compared with.
\end{itemize}\medskip
\ifdefined\hypertarget\hypertarget{sourcepointer-1}{}\fi
\noindent\textit{Source reference.} Original definition 011O: sources/stacks-homology/pinned-original-source.tex, source lines 5436--5436. Navigation pointer invoked at proof line 6074; the counted span contains the pointer only and the cited result is not reproduced inside it, so the pointer is registered as a non-load-bearing navigation stub in the established form for Stacks material. No bibliography citation is involved.\par\medskip
\ifdefined\hypertarget\hypertarget{sourcepointer-2}{}\fi
\noindent\textit{Source reference.} Original lemma 012C: sources/stacks-homology/pinned-original-source.tex, source lines 5966--5966. Navigation pointer invoked at proof line 6034; the counted span contains the pointer only and the cited result is not reproduced inside it, so the pointer is registered as a non-load-bearing navigation stub in the established form for Stacks material. No bibliography citation is involved.\par\medskip
\ifdefined\hypertarget\hypertarget{sourcepointer-3}{}\fi
\noindent\textit{Source reference.} Original definition 012E: sources/stacks-homology/pinned-original-source.tex, source lines 6038--6038. Navigation pointer invoked at proof line 6118; the counted span contains the pointer only and the cited result is not reproduced inside it, so the pointer is registered as a non-load-bearing navigation stub in the established form for Stacks material. No bibliography citation is involved.\par\medskip
\begin{lemma}
\label{lemma-compute-filtered-cohomology}
Let $\mathcal{A}$ be an abelian category. Let $(K, F, d)$ be a filtered
differential object of $\mathcal{A}$. If $Z_\infty^p$ and $B_\infty^p$
exist (see proof), then
\begin{enumerate}
\item the limit $E_\infty$ exists and is graded having
$E_\infty^p = Z_\infty^p/B_\infty^p$ in degree $p$, and
\item the associated graded $\text{gr}(H(K))$ of the cohomology of $K$
is a graded subquotient of the graded limit object $E_\infty$.
\end{enumerate}
\end{lemma}

\begin{proof}
The objects $Z_\infty$, $B_\infty$, and the limit
$E_\infty = Z_\infty/B_\infty$ of
Definition \ref{definition-limit-spectral-sequence}
are objects of $\text{Gr}(\mathcal{A})$ by our construction of
the spectral sequence in the proof of
Lemma \ref{lemma-spectral-sequence-filtered-differential}.
Since $Z_r^p$, resp.\ $B_r^p$ is the $p$th graded part of $Z_r$, resp.\  $B_r$,
if we assume that
$$
Z_\infty^p = \bigcap\nolimits_r Z_r^p =
\frac{\bigcap_r (F^pK \cap d^{-1}(F^{p + r}K) + F^{p + 1}K)}{F^{p + 1}K}
$$
and
$$
B_\infty^p = \bigcup\nolimits_r B_r^p =
\frac{\bigcup_r (F^pK \cap d(F^{p - r + 1}K) + F^{p + 1}K)}{F^{p + 1}K}.
$$
exist, then $Z_\infty$ and $B_\infty$ exist with degree $p$ parts
$Z_\infty^p$ and $B_\infty^p$ (follows from an elementary argument
about unions and intersections of graded subobjects). Thus
$$
E_\infty^p =
\frac{\bigcap_r (F^pK \cap d^{-1}(F^{p + r}K) + F^{p + 1}K)}
{\bigcup_r (F^pK \cap d(F^{p - r + 1}K) + F^{p + 1}K)}.
$$
where the top and bottom exist. We have
\begin{equation}
\label{equation-on-top}
\Ker(d) \cap F^pK + F^{p + 1}K
\subset
\bigcap\nolimits_r \left(F^pK \cap d^{-1}(F^{p + r}K) + F^{p + 1}K\right)
\end{equation}
and
\begin{equation}
\label{equation-at-bottom}
\bigcup\nolimits_r \left(F^pK \cap d(F^{p - r + 1}K) + F^{p + 1}K\right)
\subset
\Im(d) \cap F^pK + F^{p + 1}K.
\end{equation}
Thus a subquotient of $E_\infty^p$ is
$$
\frac{\Ker(d) \cap F^pK + F^{p + 1}K}{\Im(d) \cap F^pK + F^{p + 1}K} =
\frac{\Ker(d) \cap F^pK}{\Im(d) \cap F^pK + \Ker(d) \cap F^{p + 1}K}
$$
Comparing with the formula given for $\text{gr}^pH(K)$ in the discussion
following
Definition \ref{definition-filtration-cohomology-filtered-differential}
we conclude.
\end{proof}

\end{document}
